\documentclass[
  12pt,
  a4paper,
  oneside,
  brazil,
  chapter=TITLE
]{abntex2}

\usepackage[T1]{fontenc}
\usepackage[utf8]{inputenc}
\usepackage{newtxtext, newtxmath}
\usepackage[brazil]{babel}
\usepackage{microtype}
\usepackage{graphicx}
\usepackage{amsmath,amssymb}
\usepackage{booktabs}
\usepackage{indentfirst}
\usepackage{float}
\usepackage[num]{abntex2cite}
\usepackage[top=3cm, bottom=2cm, left=3cm, right=2cm]{geometry}
\usepackage{fancyhdr}
\usepackage{hyperref}

\renewcommand{\ABNTEXchapterfont}{\rmfamily}
\renewcommand{\ABNTEXchapterfontsize}{\normalsize}
\renewcommand{\ABNTEXsectionfont}{\rmfamily}
\renewcommand{\ABNTEXsectionfontsize}{\normalsize}
\renewcommand{\ABNTEXsubsectionfont}{\rmfamily}
\renewcommand{\ABNTEXsubsectionfontsize}{\normalsize}
\renewcommand{\ABNTEXsubsubsectionfont}{\rmfamily}
\renewcommand{\ABNTEXsubsubsectionfontsize}{\normalsize}

\fancypagestyle{abntpaginas}{%
  \fancyhf{}
  \fancyhead[R]{\thepage}
  \renewcommand{\headrulewidth}{0pt}
}

\hypersetup{
  colorlinks=true,
  linkcolor=black,
  citecolor=black,
  urlcolor=black
}

\setlength{\parindent}{1.25cm}
\setlength{\parskip}{0pt}

\titulo{Título do Trabalho}
\autor{Nome do Autor}
\instituicao{Nome da Instituição}
\local{Cidade}
\data{2025}
\tipotrabalho{Trabalho Avaliativo}
\preambulo{%
  Trabalho avaliativo apresentado à disciplina de
  Nome da Disciplina do curso de Nome do Curso.%
}

\renewcommand{\imprimircapa}{%
  \begin{capa}
    \centering
    {\bfseries\large\MakeUppercase{\imprimirinstituicao}\par}
    {\large Nome do Departamento\par}
    \vspace*{\fill}
    {\large\MakeUppercase{\imprimirautor}\par}
    \vspace*{\fill}
    {\bfseries\large\MakeUppercase{\imprimirtitulo}\par}
    \vspace*{\fill}
    {\large\MakeUppercase{\imprimirlocal}\par}
    {\large\imprimirdata\par}
  \end{capa}
}

\renewcommand{\imprimirfolhaderosto}{%
  \begin{folhaderosto}
    \centering
    {\large\MakeUppercase{\imprimirautor}\par}
    \vspace*{\fill}
    {\bfseries\large\MakeUppercase{\imprimirtitulo}\par}
    \vspace*{\fill}
    {\noindent\hspace{.45\textwidth}%
    \begin{minipage}{.55\textwidth}
      \SingleSpacing\small\imprimirpreambulo
    \end{minipage}\par}
    \vspace*{\fill}
    {\large\MakeUppercase{\imprimirlocal}\par}
    {\large\imprimirdata\par}
  \end{folhaderosto}
}

\begin{document}

\selectlanguage{brazil}
\OnehalfSpacing
\frenchspacing

\pretextual

\imprimircapa
\imprimirfolhaderosto

\begin{resumo}
\SingleSpacing
Texto do resumo em parágrafo único. Apresentar objetivo,
método, resultados e conclusões em voz ativa, terceira pessoa,
sem citações. De 150 a 250 palavras.

\vspace{\onelineskip}
\noindent
\textbf{Palavras-chave}: palavra-chave 1. Palavra-chave 2. Palavra-chave 3.
\end{resumo}

\pdfbookmark[0]{\contentsname}{toc}
\tableofcontents*
\cleardoublepage

\textual
\pagestyle{abntpaginas}

\chapter{Introdução}
\label{cap:introducao} 
\label{cap:introducao} 

A introdução delimita o tema, apresenta a justificativa, os
objetivos e a organização do texto \cite{ABNT14724:2011}.

\section{Objetivo Geral}

Descrever o propósito central do trabalho com um verbo no
infinitivo (analisar, desenvolver, avaliar, propor etc.).

\section{Objetivos Específicos}

\begin{itemize}
  \item Primeiro objetivo específico;
  \item Segundo objetivo específico;
  \item Terceiro objetivo específico.
\end{itemize}

\section{Justificativa}

Apresentar as razões pelas quais o tema é relevante do ponto
de vista científico, social e/ou econômico.

\chapter{Desenvolvimento}
\label{cap:desenvolvimento}

Desenvolver o conteúdo do trabalho. Esta seção pode ser
dividida em subseções conforme necessário.

Citação indireta: segundo \cite{GIL2010}, descrever
a ideia do autor com suas próprias palavras.

A \textbf{citação direta curta} (até 3 linhas) fica entre
aspas no corpo do texto: ``Exemplo de citação
direta.'' \cite[p.~10]{MARCONI2017}.

A \textbf{citação direta longa} (mais de 3 linhas) recebe
recuo de 4~cm, fonte tamanho~10 e espaço simples, sem aspas:

\begin{citacao}
  Exemplo de citação direta longa. Trecho recuado conforme
  ABNT NBR 10520:2023, com fonte menor e espaçamento simples,
  sem uso de aspas \cite[p.~21]{MINAYO2014}.
\end{citacao}

\chapter{Conclusão}
\label{cap:conclusao}

Retomar os objetivos propostos e verificar em que medida foram
atingidos. Não introduzir informações novas, apenas sintetizar
as principais descobertas do trabalho.

\postextual

\bibliography{referencias}

\end{document}